\chapter{Aiming the engine}

\textbf{Purpose:} Separate the engine's capability from knowing where to point it, and record what one night of pointing it badly cost and bought. \textbf{Scope:} the arbitrary-precision engine, and the ten rules the session's six faults produced. Nothing here is about hashing.

The engine computes at a width nothing on a desktop is supposed to reach. That capability is separate from knowing where to point it, and one night of pointing it badly produced six findings that were all faults in the pointing.

\section{What the engine is, stated so the rest has a subject}

Arbitrary-precision integer arithmetic with no floor under it, a device multiply picking its own arm by operand size, and division done by Newton so that multiplies and shifts are the only primitives anywhere. Measured on one desktop:

\begin{center}
\begin{tabular}{@{}ll@{}}
\toprule
what & measured \\
\midrule
pi to \(10{,}000{,}000\) digits & \(116.5\) s, checked twice against an independent series \\
a \(1{,}290{,}000{,}000\)-bit product & \(1.72\) s of device time \\
\(2.15\) trillion SHA-256 evaluations & summarized in \(3\) MB of counters \\
\(26.6\) billion multiply-accumulates & \(3.2\) s, once the work was in the shape the card wants \\
\bottomrule
\end{tabular}
\end{center}

The last line rewards reading twice. The same permutation null took \(24\) draws a minute in Python and \(2000\) draws in three seconds as a Gram matrix. The card was never the constraint; the shape of the work was.

\section{Read the interior, not the boundary}

A physical system forces a boundary reading: the interior cannot be reached and is inferred from what escapes. Almost nothing the engine is pointed at has that constraint, and treating it as though it does is a blindfold put on by hand.

Pointed at a compression function, hours went into reading digests while every schedule word, every state word at every round, and every value built inside a round were readable at every instant. The interior said in one run what the boundary could not say at all: the two nonlinear functions lag the linear ones by exactly one round, because a value like \((e \wedge f) \oplus (\neg e \wedge g)\) does not depend on \(e\) wherever \(f\) and \(g\) agree. A digest sums every path into one number and can never separate them.

\textbf{Ask what the object's interior is before deciding to read its surface.}

\section{Ask what is forbidden before asking what is frequent}

A frequency needs a null, a bar, an argument and a correction. A support does not: a transition either occurred or it never did, and a transition that cannot occur rules out every trajectory requiring it, permanently, with no statistics involved.

Six hundred excitations found \(214\) of \(256\) positions unreachable at one depth. That is a harder statement than any \(p\)-value produced all night, and it took one run.

\textbf{Frequencies are what is left when the rules run out. Read the rules first.}

\section{The positive control decides whether a null means anything}

Without one, \emph{we found nothing} and \emph{we could not have found anything} print identically, and there is no way to tell them apart from inside the result.

The clearest case: a compression test found no structure in a hash and was about to be believed. Adding an arm for the digits of pi, a sequence with a \emph{kilobyte} generating program by construction, showed pi compressing to \(1.0010\) against random's \(1.0001\). The instrument could not distinguish a known short program from none. Its verdict was empty, and the claim would have been a fact about zlib.

Every instrument gets one: pi for algorithmic structure, a planted relation for a relation search, a held-out range for anything selected, a known-answer vector for any reader, planted structure of each class for a scanner. If the control does not fire, the null is about the tool.

\section{Draw the bar, never derive it}

Seven derived thresholds in one session and every one came in too low. The error runs in one direction only, and that earns a rule instead of more care. Deriving a bar means enumerating the sources of variance, and the ones left out only ever \emph{add} variance. The derived value is always the low estimate. A drawn bar has them all whether or not anyone thought of them.

\begin{center}
\begin{tabular}{@{}l@{}}
\toprule
the seven, by shape \\
\midrule
a shuffle of labels returning identical counts every draw, sd \(0.00\) \\
a bar at two Poisson floors for a statistic spanning three and a half by construction \\
chi-square's variance through a logarithm: \(0.0582\) where shuffling gives \(0.1015\) \\
\(\sqrt{2 \ln N}\), the \emph{expected} maximum, which half of all nulls exceed \\
\bottomrule
\end{tabular}
\end{center}

Resample instead. Shuffle the labels, redraw the counts at the same \(N\) and the same bin count, recompute the identical statistic, read the percentile. Where the statistic is a maximum, a minimum or a range this is not optional, because those have no usable closed form.

\section{Match the null to the nuisance parameters}

A null that does not match reports the nuisance parameter and calls it structure. A lit set's spectrum depends heavily on how many bits are lit, and the weight wandered between \(57\) and \(140\) across the rounds being compared. A null drawn at a fixed weight would therefore have reported that wander as signal at every step.

\textbf{Hold fixed whatever the statistic is sensitive to and does not mean.}

\section{Normalize until the baseline is flat}

For points with no preferred direction every mode carries the same expected power, and a degree holds \(2l+1\) of them. Per-degree power therefore \emph{rises} as \(2l+1\) with nothing in it at all. Drawn: the ratio is flat at \(3.867\) for a weight of \(64\) and \(5.095\) for \(128\), to within two and a half percent across ten degrees, and the two constants stand in the ratio the finite-population correction predicts, \(1.318\) measured against \(1.333\) expected.

Plotted raw, every spectrum rises and the rise is entirely mode count. Divided by \(2l+1\) the null is a horizontal line, the baseline is one number instead of a curve, and a departure is a departure with nothing to subtract by eye.

\textbf{Find the form in which the null is flat, then plot that.}

\section{Know which class an instrument can see}

Scanners were raced against planted structure of four kinds. The harmonic expansion beat every rival on anything with a shape, \(328\) standard errors against \(14.8\) for the next best, and on a forced parity it read \(2.43\) against its own \(3.75\) floor, \emph{below} its noise. Not worse: blind, because the expansion is a linear map of the lit set and a parity is not a linear property of which bits are lit. The scan that caught it was the crudest in the table.

\emph{Best instrument} is therefore not a property a single transform has. Run one cheap and complete scan over the spatial class and a second over the algebraic class, and report which fired.

\section{Distinguish the format's floor from the object's}

A residual \emph{stopping} at \(10^{-16}\) is float64. A residual \emph{falling} with the width belongs to the object. Only fixed point can tell those apart, and the difference is not academic: an invariance measured in a double reported \(10^{-14}\) and could not say whether the theorem held exactly or merely sat where a double runs out. Recomputed in fixed point it read \(2^{-60}\), \(2^{-92}\) and \(2^{-123}\) at widths of \(64\), \(96\) and \(128\), tracking the width. The invariance is exact.

\textbf{Measure the same thing at two widths. If the answer moves, it was the format.}

\section{Widen everything that grows, and nothing that is defined}

A mod \(2^{32}\) ring \emph{is} the function; widening it computes a different one and fails silently. But the statistics computed \emph{over} that ring grow without bound and must not share its width. An accumulator held in a thirty-two bit integer was nine bits from silently wrapping, in analysis code, while the argument for keeping the ring narrow was being made about the same number.

\textbf{The definition fixes a width. Everything else is arbitrary precision.}

\section{The instrument is the first suspect}

Six faults in one session, every one producing something that looked like a finding, every one caught by the result being too clean and never by care beforehand:

\begin{center}
\begin{tabular}{@{}l@{}}
\toprule
what it looked like, and what it was \\
\midrule
a thirty sigma excess that was a generator with a period of \(256\) \\
twenty hours of geography that were twenty random trough hours \\
a swing of \(1.394\) that was the null's own largest \\
a band at \(280\) seconds that was the top of a sorted list \\
a chirality result fixed by geometry before the data was read \\
an invariance that was a scalar hiding the structure it summed over \\
\bottomrule
\end{tabular}
\end{center}

When a result is surprising, the tool produced it and the tool is what changed. Check it first, and build the check into the tool so the next surprise has to pass it.

\section{A memoryless corpus, and what that buys}

A corpus is memoryless when its items carry nothing about each other: drawing one tells you nothing about the next. Establishing it is worth doing early, because it decides which whole families of method are available and which are wasted effort.

It was established here by direct measurement instead of assumption. Distance-to-target autocorrelation across \(200{,}000\) consecutive items was inside the null band at every lag from \(1\) to \(1024\). Hill-climbing on that distance lost to random draw by \(3.81\) standard errors on an equal budget. Not tied, \emph{lost}, because it spent its budget probing neighbors that carried nothing.

\subsection{Ruled out, at and above a measured floor}

Every ordering strategy at once. If items are independent, search order cannot change the expected cost, and any cleverness about \emph{which} item to try next is spending effort on a quantity that does not exist. One measurement retires the entire family instead of each member separately, the strongest thing a null does anywhere in this tree, and the reason it got quoted past its evidence.

Both arms read the digest through the leading-zero count, about two bits of entropy per evaluation, discarding everything below the first set bit. That was never wrong, and it was never bounded either: \emph{we found no gradient} was standing in for \emph{no gradient exists} with nothing fixing the difference. Fading an injected gradient until the same statistic loses it puts a number on it: \(0.0625\) leading-zero bits per single-bit step, with the control flat at every rung. Above that, excluded. Below it, or expressed anywhere else in the digest, not looked at.

The engineering conclusion survives intact, because mining's objective \emph{is} that count. The general claim did not, and the difference carries the lesson: this null was the result carrying the most weight here and the last one to get a floor, for the reason every fault in this chapter shares. It looked clean, and a readout guaranteeing a near-null looks the same as an object with nothing in it.

\subsection{Granted: the arms}

Independence lets observations be pooled by addition with no calibration step. A memoryless corpus can therefore be sampled in any order, at any depth, by any number of observers, and combined exactly.

This half is the useful one and it is easy to miss. Because each arm carries its own exact zero, a count \(c\) out of \(N\) deviating by exactly \(2c - N\), arms of \emph{different} depths pool correctly, the deviations add, the variances add with them, and there is no fitting anywhere. Nothing has to be planned in advance: short arms while exploring, long arms where something looks interesting, all combined into one exact estimate.

Memory per arm is flat. The counters are the same size whatever the depth, and a thousand arms therefore hold \(2.15\) trillion evaluations in three megabytes, about \(700\) million per kilobyte, and the footprint does not grow. That makes \emph{keep looking until the question is answered} a real option instead of a wish: the record never fills.

And the two things arms buy do not scale alike, which decides how many to run:

\begin{center}
\begin{tabular}{@{}ll@{}}
\toprule
what it buys & how it scales \\
\midrule
pooled depth & as \(\sqrt{k}\), the slow axis, traversed at whatever rate the device runs \\
agreement & as \(2^{k}\), each further arm halving the chance a run of agreeing signs is coincidence \\
\bottomrule
\end{tabular}
\end{center}

The second saturates. At \(32\) arms a unanimous reading happens by chance about once in ten million over \(256\) cells, and past that point more arms buy depth only. Thirty-two is the knee.

\section{Elision, and how to use it instead of fighting it}

Elision is a value ceasing to depend on an input, the state becoming identical whether that input was one thing or another. It is normally described as loss, and from the inversion side it is. From the measurement side it is the most useful thing in the machine, because an elided dependency is a \emph{hard} fact about the object and costs nothing to establish.

The support measurement is elision read forwards. Six hundred excitations, and \(214\) of \(256\) positions never moved at all, each one an elision, each one ruling out every trajectory that would require it, and none of them needing a null or a bar or a \(p\)-value. A frequency needs all three and most come back empty anyway.

Three ways to use it:

\textbf{As a constraint, and not a loss.} Where a dependency has been elided, everything downstream of that dependency is decided. That prunes searches over trajectories outright instead of probabilistically, and it makes a support measurement worth more per run than any frequency.

\textbf{As a clock.} Elision accumulates at a rate, and the rate is a property worth knowing. Here the support grew about nine positions a round and closed near round thirty of sixty-four, saying the object spends half its work reaching the point where everything can affect everything and the rest past it. That number is invisible to any reading of the finished output.

\textbf{As the thing that tells you where to stop measuring.} Once a dependency is elided no instrument recovers it, however precise, because the state is genuinely identical either way. A measurement aimed past the elision point is measuring nothing and will return a floor that looks like a result. Knowing where elision completes is knowing where the readable window closes, and in the case measured here it closed at round thirty while the tempting place to look was round sixty-four.

The general form: \textbf{elision is where the object stops being able to tell you something, and finding that boundary is cheaper and harder-edged than anything on either side of it.}

\section{The order these go in}

\begin{enumerate}
  \item State what the object permits before measuring what it does.
  \item Read the interior if the interior is readable.
  \item Run the positive control. If it does not fire, stop.
  \item Draw the bar. Match the null to the nuisance parameters.
  \item Normalize until the baseline is flat.
  \item Measure at two widths and see whether the floor moves.
  \item When something clears, suspect the instrument, and only then the object.
\end{enumerate}
